\chapter{Analyse und Reverse Engineering}\label{ch:analyse-und-reverse-engineering}
Ziel der Analyse ist es, die auf der Platine vorhandenen Bauteile und deren elektrische Verbindungen systematisch zu erfassen.
Hierzu werden die Leiterbahnen schrittweise nachverfolgt und die Ergebnisse mithilfe der Erfassungs-App dokumentiert.
Dabei entstehen eine Bauteilliste und eine Verbindungsliste, die als Grundlage für die weitere Rekonstruktion der Schaltung dienen.

\section{Demontage und Vorbereitung der Platine}\label{sec:demontage-und-vorbereitung-der-platine}
Zunächst wurden die Pin-Stecker vorsichtig von der Platine getrennt.
Anschließend wurden die vier Schrauben an den Ecken gelöst und zusammen mit den beiden Gewindeelementen zur Höhenanpassung geordnet neben dem Modell abgelegt.
Danach konnte die Platine vom Modell getrennt und für die Untersuchung zugänglich gemacht werden.

\section{Vorgehensweise bei der Bauteil- und Verbindungserfassung}\label{sec:vorgehensweise-bei-der-bauteil--und-verbindungserfassung}
Die Aufnahme erfolgte arbeitsteilig durch zwei Gruppenmitglieder.
Während eine Person die Leiterbahnen auf der Platine nachverfolgte und die angeschlossenen Bauteile identifizierte, dokumentierte die zweite Person die Ergebnisse in der App.

Hierfür wurde eine Leiterbahn jeweils so weit verfolgt, bis ein noch nicht erfasstes Bauteil erreicht wurde.
Dieses wurde beschrieben, mit einer eindeutigen ID versehen und in die Bauteilliste aufgenommen.
Anschließend konnten die miteinander verbundenen Bauteile in der App ausgewählt und ihre Verbindung dokumentiert werden.
Soweit erforderlich, wurde dabei der jeweilige Anschlusspin berücksichtigt.

\section{Untersuchung der Spannungsversorgung}\label{sec:untersuchung-der-spannungsversorgung}
Die Untersuchung begann mit der Leiterbahn der −12-V-Spannungsversorgung.
Ausgehend vom Versorgungsanschluss wurden die Leiterbahnen und alle unmittelbar damit verbundenen Bauteile systematisch erfasst.
Dieses Vorgehen wurde anschließend für die weitere Versorgungsspannung und die Masseverbindungen wiederholt.

Dadurch konnte bereits ein großer Teil der Bauteile aufgenommen und deren Anbindung an die Spannungsversorgung beziehungsweise an Masse dokumentiert werden.
Gleichzeitig entstand eine erste Orientierung über den Aufbau der Platine.

\section{Untersuchung der Pin-Stecker}\label{sec:untersuchung-der-pin-stecker}
Im nächsten Schritt wurden die Anschlüsse der Pin-Stecker untersucht.
Für jeden Anschluss wurde nachverfolgt, mit welchen Bauteilen er elektrisch verbunden ist.
Die Erfassung erfolgte nach demselben Prinzip wie bei der Untersuchung der Spannungsversorgung.

Bereits dokumentierte Bauteile wurden anhand ihrer IDs zugeordnet.
Neu gefundene Bauteile wurden ergänzt und ihre Verbindungen eingetragen.
Mit fortschreitender Aufnahme verlagerte sich der Schwerpunkt zunehmend auf die Nachverfolgung der Verbindungen zwischen den bereits bekannten Bauteilen.

\section{Fotografische Dokumentation und Bauteilzuordnung}\label{sec:fotografische-dokumentation-und-bauteilzuordnung}
Nach der Untersuchung der Pin-Stecker wurde ein Foto der Platinenoberseite aufgenommen.
Auf diesem wurden die bereits erfassten Bauteile mit ihren zugewiesenen IDs beschriftet (siehe Abbildung zur Bauteilzuordnung).

Die beschriftete Aufnahme erleichterte die Orientierung auf der Platine und den Abgleich mit der Bauteilliste.
Dabei konnten zuvor übersehene Bauteile erkannt und in der Dokumentation ergänzt werden.

\section{Untersuchung der integrierten Schaltungen}\label{sec:untersuchung-der-integrierten-schaltungen}
Anschließend wurden die integrierten Schaltungen (ICs) einzeln untersucht.
Für jeden IC wurden die Verbindungen an den jeweiligen Anschlusspins nachverfolgt und in der App dokumentiert.

Die dabei gewonnenen Zuordnungen wurden außerdem in der beschrifteten Platinenaufnahme ergänzt.
Dadurch ließ sich die Übersicht über die Verbindungen zwischen den ICs und den übrigen Bauteilen schrittweise erweitern.

\section{Zusammenführung und Prüfung der erfassten Daten}\label{sec:zusammenfuhrung-und-prufung-der-erfassten-daten}
Die Bauteilliste, die Verbindungsliste und die beschriftete Platinenaufnahme bilden gemeinsam die Grundlage für die weitere Auswertung.
Zur abschließenden Prüfung werden diese Unterlagen miteinander abgeglichen.
Dabei ist insbesondere zu kontrollieren, ob alle Bauteile eindeutig zugeordnet und die aufgenommenen Verbindungen vollständig und widerspruchsfrei dokumentiert sind.

Noch unklare Verbindungen und nicht eindeutig identifizierte Bauteile werden gekennzeichnet und gezielt nachuntersucht.
Auf dieser Grundlage kann anschließend die Schaltungsstruktur weiter rekonstruiert werden.

% Zum Einfügen in die bestehende Dokumentation mit \usepackage{settings}.
% Alle 168 IDs aus der CSV bleiben erhalten. Stand: 27.10.2026.
\subsection{Bauteilliste der untersuchten Platine}
\label{sec:bauteilliste}

Die folgende Liste enthält alle 168 in der Aufnahme vergebenen IDs.
Jede ID wurde im Export mit der Anzahl 1 erfasst. Auch Anschlüsse,
Steckverbinder und Drahtbrücken bleiben für den Abgleich mit der
Verbindungsliste enthalten. Die Zahl der IDs ist daher nicht mit
der Anzahl eigenständiger elektronischer Bauteile gleichzusetzen.

Die Werte wurden aus den notierten Aufdrucken und Farbringen abgeleitet;
eine messtechnische Bestätigung liegt nicht vor. Widerstände mit vier
Ringen werden mit zwei, solche mit fünf Ringen mit drei Wertziffern
ausgewertet. Bei sechs Ringen werden die ersten fünf entsprechend
ausgewertet; die Bedeutung des Zusatzrings ist noch zu prüfen.
Die Tabelle nennt allgemeine Bauteilfunktionen, soweit sie durch den Typ
belegt sind. Die tatsächliche Funktion in der Platine folgt erst aus
der Untersuchung der Beschaltung.

\textbf{Kennzeichnung:} \emph{Farbcode} = aus der Abschrift berechneter
Nennwert; \emph{Aufdruck} = Wert aus der Aufnahme;
\emph{Typ} = allgemeine Bauteilart durch Datenblatt zugeordnet;
\emph{vorläufig} = plausible, noch zu bestätigende Zuordnung;
\emph{offen} = keine eindeutige Bestimmung möglich.
Q1--Q17 verweisen auf die Quellen unterhalb der Tabelle.

\begin{GPMTable}
{L{1.0cm}|L{3.0cm}|L{3.0cm}|Y|L{1.0cm}}
{Bauteile, Anschlüsse und Verbindungselemente der untersuchten Platine}
{tab:bauteilliste}
{\TableHead{ID} & \TableHead{Bauteilart} & \TableHead{Wert / Typ} & \TableHead{Zuordnung und Prüfhinweise} & \TableHead{Quelle}}
    \textbf{P001} & Sicherung & Nennstrom unbekannt & offen. Bauform, Auslösecharakteristik und Nennstrom ablesen. & -- \\
    \hline
    \textbf{P002} & Elektrolytkondensator & 220 \ensuremath{\mu}F / 40 V & Aufdruck. Polung und tatsächliche Versorgungsknoten prüfen. Die CSV nennt +12 V und \ensuremath{-}12 V. & -- \\
    \hline
    \textbf{P003} & Schaltregler-IC & TL497ACN & Typ. Für Aufwärts-, Abwärts- oder invertierende Wandler; konkrete Schaltung noch bestimmen. & Q2 \\
    \hline
    \textbf{P004} & Kondensator, vermutlich & 470 nF & vorläufig. Aufdruck 470N als Kapazität interpretiert; Bauart und Spannungsfestigkeit offen. & -- \\
    \hline
    \textbf{P005} & Widerstand & vermutlich 4,7 k\ensuremath{\Omega} & vorläufig. Farbfolge enthält „leer“. Bei fünf relevanten Ringen Gelb–Lila–Schwarz–Braun–Braun: 4,7 k\ensuremath{\Omega} \ensuremath{\pm}1 \%. & Q1 \\
    \hline
    \textbf{P006} & Elektrolytkondensator & 1000 \ensuremath{\mu}F / 6,3 V & Aufdruck. Temperaturangabe laut Aufnahme: 85 °C. & -- \\
    \hline
    \textbf{P007} & Schottkydiode & BAT41 & Typ. In der CSV als Freilaufdiode bezeichnet; diese Schaltungsfunktion ist noch zu prüfen. & Q9 \\
    \hline
    \textbf{P008} & Widerstand & 10 k\ensuremath{\Omega}; \ensuremath{\pm}1 \% & Farbcode.  Ringe: Braun-Schwarz-Schwarz-Rot-Braun. & Q1 \\
    \hline
    \textbf{P009} & Versorgungsanschluss & +12 V & Aufdruck. Erfasster Anschluss, kein eigenständiger Halbleiter. & -- \\
    \hline
    \textbf{P010} & Versorgungsanschluss & \ensuremath{-}12 V & Aufdruck. Bedeutung als negative Schiene oder Versorgungsrückleiter prüfen. & -- \\
    \hline
    \textbf{P011} & Draht / Drahtbrücke & blau & Aufdruck. Leiterquerschnitt und Funktion nicht dokumentiert. & -- \\
    \hline
    \textbf{P012} & Masseanschluss & Analogausgang & Aufdruck. Als „Erde“ erfasst; Schutzleiterfunktion nicht nachgewiesen. & -- \\
    \hline
    \textbf{P013} & Widerstand & 220 k\ensuremath{\Omega}; \ensuremath{\pm}1 \% & Farbcode.  Ringe: Rot-Rot-Schwarz-Ornage-Braun. & Q1 \\
    \hline
    \textbf{P014} & Spannungsreferenz & LM329DZ / ca. 6,9 V & Typ. Kein diskreter Transistor. „524“ ist Zusatzaufdruck, kein Referenzspannungswert. & Q3 \\
    \hline
    \textbf{P015} & Widerstand & 100 k\ensuremath{\Omega}; \ensuremath{\pm}1 \% & Farbcode.  Ringe: Braun-Schwarz-Schwarz-Orange-Braun. & Q1 \\
    \hline
    \textbf{P016} & Diode, Typ ungeklärt & möglicherweise 1N4148 & vorläufig. „PN4841“ und ähnliche Lesarten könnten zum Aufdruck 4148PH gehören. Nicht als Z-Diode bestätigt. Aufdruck/Foto prüfen. & Q10 \\
    \hline
    \textbf{P017} & IC, vermutlich OPV & möglicherweise TAE1453A & vorläufig. Erfasst als TAB1453A. TAE1453 ist ein OPV mit offenem Kollektorausgang; genaue Beschriftung und Pinzahl prüfen. & Q17 \\
    \hline
    \textbf{P018} & Logik-IC & CD4001B & Typ. Vier NOR-Gatter mit je zwei Eingängen; Beschaltung als Inverter ist möglich. Zusatzaufdrucke nicht als Typ übernommen. & Q4 \\
    \hline
    \textbf{P019} & IC, vermutlich OPV & möglicherweise TAE1453A & vorläufig. Erfasst als TAB1453A. TAE1453 ist ein OPV mit offenem Kollektorausgang; genaue Beschriftung und Pinzahl prüfen. & Q17 \\
    \hline
    \textbf{P020} & Leistungs-MOSFET, vermutlich & BUZ10 & vorläufig. „BUZ1U“ bzw. „BUZ 108048“ vermutlich BUZ10 mit Zusatzcode 8048. Vollständigen Aufdruck prüfen. & Q13 \\
    \hline
    \textbf{P021} & Spannungsreferenz & LM329DZ / ca. 6,9 V & Typ. Kein diskreter Transistor. „524“ ist Zusatzaufdruck, kein Referenzspannungswert. & Q3 \\
    \hline
    \textbf{P022} & Widerstand & 10 k\ensuremath{\Omega}; \ensuremath{\pm}1 \% & Farbcode.  Ringe: Braun-Schwarz-Schwarz-Rot-Braun. & Q1 \\
    \hline
    \textbf{P023} & Widerstand & 430 k\ensuremath{\Omega} & vorläufig. Bei Silber: \ensuremath{\pm}10 \%; bei Grau: \ensuremath{\pm}0,05 \%. Farbring und Leserichtung prüfen. & Q1 \\
    \hline
    \textbf{P024} & Zweifacher CMOS-OPV & ICL7621 & Typ. „I“ in der Aufnahme vermutlich Ziffer 1; vollständige Endung des Typs prüfen. & Q5 \\
    \hline
    \textbf{P025} & Steckverbinder & weiß & Aufdruck. Laut Aufnahme sind Pin 19 und 20 unbelegt; Pinzahl und Belegung prüfen. & -- \\
    \hline
    \textbf{P026} & Leistungs-MOSFET, vermutlich & BUZ10A & vorläufig. BUZ10A deutet auf Einzel-MOSFET. Die Bezeichnung „Transistoren-Pack“ und Stückzahl am Bauteil prüfen. & Q13 \\
    \hline
    \textbf{P027} & Logik-IC & CD4001B & Typ. Vier NOR-Gatter mit je zwei Eingängen; Beschaltung als Inverter ist möglich. Zusatzaufdrucke nicht als Typ übernommen. & Q4 \\
    \hline
    \textbf{P028} & Steckverbinder & mehrfarbiges Kabel & Aufdruck. Pinzahl, Rastermaß und Belegung offen. & -- \\
    \hline
    \textbf{P029} & Logik-IC & CD4001B & Typ. Vier NOR-Gatter mit je zwei Eingängen; Beschaltung als Inverter ist möglich. Zusatzaufdrucke nicht als Typ übernommen. & Q4 \\
    \hline
    \textbf{P030} & N-Kanal-MOSFET & BSS89 & Typ. Selbstsperrender Kleinsignal-MOSFET; Pinbelegung herstellerabhängig prüfen. & Q12 \\
    \hline
    \textbf{P031} & Kondensator, Bauart offen & 1 \ensuremath{\mu}F & Aufdruck. 1u0 = 1 \ensuremath{\mu}F; zusätzlich „40VR, 25VCN“ notiert. Spannungsangaben und Bauart prüfen. & -- \\
    \hline
    \textbf{P032} & Spannungsregler, vermutlich & 78Lxx, eventuell 78L05 & vorläufig. Aufnahme „78L 0,5A“ unvollständig. Ein 78L05 wäre ein 5-V-/100-mA-Regler; 0,5 A ist dadurch nicht belegt. & Q16 \\
    \hline
    \textbf{P033} & Kondensator, Bauart offen & 1 \ensuremath{\mu}F & Aufdruck. 1u0 = 1 \ensuremath{\mu}F; zusätzlich „40VR, 25VCN“ notiert. Spannungsangaben und Bauart prüfen. & -- \\
    \hline
    \textbf{P034} & Kondensator, Bauart offen & 1 \ensuremath{\mu}F & Aufdruck. 1u0 = 1 \ensuremath{\mu}F; zusätzlich „40VR, 25VCN“ notiert. Spannungsangaben und Bauart prüfen. & -- \\
    \hline
    \textbf{P035} & Kondensator, Bauart offen & 1 \ensuremath{\mu}F & Aufdruck. 1u0 = 1 \ensuremath{\mu}F; zusätzlich „40VR, 25VCN“ notiert. Spannungsangaben und Bauart prüfen. & -- \\
    \hline
    \textbf{P036} & Kondensator, Bauart offen & 1 \ensuremath{\mu}F & Aufdruck. 1u0 = 1 \ensuremath{\mu}F; zusätzlich „40VR, 25VCN“ notiert. Spannungsangaben und Bauart prüfen. & -- \\
    \hline
    \textbf{P037} & Monoflop-IC & HEF4528BP & Typ. Zwei nachtriggerbare monostabile Multivibratoren; Impulsdauer durch externe RC-Glieder. & Q6 \\
    \hline
    \textbf{P038} & Monoflop-IC & HEF4528BP & Typ. Zwei nachtriggerbare monostabile Multivibratoren; Impulsdauer durch externe RC-Glieder. & Q6 \\
    \hline
    \textbf{P039} & Monoflop-IC & HEF4528BP & Typ. Zwei nachtriggerbare monostabile Multivibratoren; Impulsdauer durch externe RC-Glieder. & Q6 \\
    \hline
    \textbf{P040} & Masseanschluss & Digitaleingang & Aufdruck. Als „Erde“ erfasst; Schutzleiterfunktion nicht nachgewiesen. & -- \\
    \hline
    \textbf{P041} & Masseanschluss & Digitalausgang & Aufdruck. Als „Erde“ erfasst; Schutzleiterfunktion nicht nachgewiesen. & -- \\
    \hline
    \textbf{P042} & Kondensator, Bauart offen & 1 \ensuremath{\mu}F / 25 V & Aufdruck. 1u0 als 1 \ensuremath{\mu}F interpretiert; Hersteller laut Aufnahme Philips. & -- \\
    \hline
    \textbf{P043} & Monoflop-IC & HEF4528BP & Typ. Zwei nachtriggerbare monostabile Multivibratoren; Impulsdauer durch externe RC-Glieder. & Q6 \\
    \hline
    \textbf{P044} & Widerstandsnetzwerk & vermutlich 3,3 k\ensuremath{\Omega} je Element & vorläufig. A332J: Code 332 = 3,3 k\ensuremath{\Omega}; J üblicherweise \ensuremath{\pm}5 \%. Anzahl, interne Verschaltung und Serie prüfen. Laut Aufnahme wenig Lötzinn an Pin 2. & Q14 Q1 \\
    \hline
    \textbf{P045} & Tantalkondensator, vermutlich & vermutlich 22 \ensuremath{\mu}F / 16 V & vorläufig. Aufdruck „+22 16V, Erotantal S0“ spricht für gepolten Tantalkondensator; Typ und Polung prüfen. & -- \\
    \hline
    \textbf{P046} & N-Kanal-MOSFET & BSS91 & Typ. Selbstsperrender Kleinsignal-MOSFET; N7 als Zusatzaufdruck erfasst. & Q12 \\
    \hline
    \textbf{P047} & N-Kanal-MOSFET & BSS91 & Typ. Selbstsperrender Kleinsignal-MOSFET; N7 als Zusatzaufdruck erfasst. & Q12 \\
    \hline
    \textbf{P048} & N-Kanal-MOSFET & BSS91 & Typ. Selbstsperrender Kleinsignal-MOSFET; N7 als Zusatzaufdruck erfasst. & Q12 \\
    \hline
    \textbf{P049} & Analogmultiplexer-IC & vermutlich HCF4051BE & vorläufig. Aufdruck HCF405IBE wahrscheinlich HCF4051BE: ein 8-Kanal-Analogmultiplexer/Demultiplexer. & Q7 \\
    \hline
    \textbf{P050} & Zweifacher CMOS-OPV & ICL7621 & Typ. „I“ in der Aufnahme vermutlich Ziffer 1; vollständige Endung des Typs prüfen. & Q5 \\
    \hline
    \textbf{P051} & Widerstand & 47 k\ensuremath{\Omega}; \ensuremath{\pm}1 \% & Farbcode.  Ringe: Gelb-Lila-Schwarz-Rot-Braun. & Q1 \\
    \hline
    \textbf{P052} & Widerstand & 1,5 k\ensuremath{\Omega}; \ensuremath{\pm}1 \% & Farbcode.  Ringe: Braun-Grün-Schwarz-Braun-Braun. & Q1 \\
    \hline
    \textbf{P053} & Operationsverstärker & TL061CP & Typ. Einfacher OPV mit JFET-Eingang; 8719CR als Zusatzaufdruck erfasst. & Q8 \\
    \hline
    \textbf{P054} & Widerstand & 10 k\ensuremath{\Omega}; \ensuremath{\pm}1 \% & Farbcode.  Ringe: Braun-Schwarz-Schwarz-Rot-Braun. & Q1 \\
    \hline
    \textbf{P055} & Widerstand & 10 k\ensuremath{\Omega}; \ensuremath{\pm}1 \% & Farbcode.  Ringe: Braun-Schwarz-Schwarz-Rot-Braun. & Q1 \\
    \hline
    \textbf{P056} & Widerstandsnetzwerk & vermutlich 33 k\ensuremath{\Omega} je Element & vorläufig. 333 = 33 × 10\textsuperscript{3} \ensuremath{\Omega}. Hersteller/Serie AM5, Anzahl und interne Verschaltung ungeklärt. & Q14 \\
    \hline
    \textbf{P057} & Widerstand & ungeklärt & offen. Notiert Orange–Orange–Schwarz–Grau/Weiß–Braun: rechnerisch 33 G\ensuremath{\Omega}/330 G\ensuremath{\Omega} \ensuremath{\pm}1 \%. Sehr ungewöhnlich; Farbringe neu ablesen. & Q1 \\
    \hline
    \textbf{P058} & Induktivität, vermutlich & Induktivität unbekannt & vorläufig. „Gewickeltes Draht Ding“ spricht für Spule/Drossel; Foto, Anschlusszahl und Messung erforderlich. & -- \\
    \hline
    \textbf{P059} & Widerstand & 120 k\ensuremath{\Omega}; \ensuremath{\pm}1 \% & Farbcode.  Ringe: Braun-Rot-Schwarz-Orange-Braun. & Q1 \\
    \hline
    \textbf{P060} & PNP-Transistor, vermutlich & BC307, Endung unklar & vorläufig. Aufdruck BC307(6) unvollständig; genaue Variante prüfen. & Q11 \\
    \hline
    \textbf{P061} & Widerstand & 27,4 k\ensuremath{\Omega}; \ensuremath{\pm}1 \% & Farbcode.  Ringe: Rot-Lila-Gelb-Rot-Braun. & Q1 \\
    \hline
    \textbf{P062} & Widerstand & 4,7 k\ensuremath{\Omega}; \ensuremath{\pm}1 \% & Farbcode.  Ringe: Gelb-Lila-Schwarz-Braun-Braun. & Q1 \\
    \hline
    \textbf{P063} & Halbleiter, nicht identifiziert & Aufdruck 452LM & offen. Die Zeichenfolge reicht zur eindeutigen Typbestimmung nicht aus; vollständigen Aufdruck erfassen. & -- \\
    \hline
    \textbf{P064} & Widerstand & 120 k\ensuremath{\Omega}; \ensuremath{\pm}1 \% & Farbcode.  Ringe: Braun-Rot-Schwarz-Orange-Braun. & Q1 \\
    \hline
    \textbf{P065} & Widerstandsnetzwerk & vermutlich 22 k\ensuremath{\Omega} je Element & vorläufig. A223J: Code 223 = 22 k\ensuremath{\Omega}; J üblicherweise \ensuremath{\pm}5 \%. Anzahl, interne Verschaltung und Serie prüfen. & Q14 Q1 \\
    \hline
    \textbf{P066} & Signalanschluss & Analogausgang U & Aufdruck. Signalpegel und Belastbarkeit noch bestimmen. & -- \\
    \hline
    \textbf{P067} & Signalanschluss & Digitaleingang 0 & Aufdruck. Signalpegel und Steckerzuordnung noch bestimmen. & -- \\
    \hline
    \textbf{P068} & Signalanschluss & Digitaleingang 1 & Aufdruck. Signalpegel und Steckerzuordnung noch bestimmen. & -- \\
    \hline
    \textbf{P069} & Signalanschluss & Digitaleingang 2 & Aufdruck. Signalpegel und Steckerzuordnung noch bestimmen. & -- \\
    \hline
    \textbf{P070} & Signalanschluss & Digitaleingang 3 & Aufdruck. Signalpegel und Steckerzuordnung noch bestimmen. & -- \\
    \hline
    \textbf{P071} & Signalanschluss & Digitaleingang 4 & Aufdruck. Signalpegel und Steckerzuordnung noch bestimmen. & -- \\
    \hline
    \textbf{P072} & Signalanschluss & Digitaleingang 5 & Aufdruck. Signalpegel und Steckerzuordnung noch bestimmen. & -- \\
    \hline
    \textbf{P073} & Signalanschluss & Digitaleingang 6 & Aufdruck. Signalpegel und Steckerzuordnung noch bestimmen. & -- \\
    \hline
    \textbf{P074} & Widerstandsnetzwerk & vermutlich 3,3 k\ensuremath{\Omega} je Element & vorläufig. A332J: Code 332 = 3,3 k\ensuremath{\Omega}; J üblicherweise \ensuremath{\pm}5 \%. Anzahl, interne Verschaltung und Serie prüfen. & Q14 Q1 \\
    \hline
    \textbf{P075} & N-Kanal-MOSFET & BSS89 & Typ. Selbstsperrender Kleinsignal-MOSFET; Pinbelegung herstellerabhängig prüfen. & Q12 \\
    \hline
    \textbf{P076} & N-Kanal-MOSFET & BSS89 & Typ. Selbstsperrender Kleinsignal-MOSFET; Pinbelegung herstellerabhängig prüfen. & Q12 \\
    \hline
    \textbf{P077} & N-Kanal-MOSFET & BSS89 & Typ. Selbstsperrender Kleinsignal-MOSFET; Pinbelegung herstellerabhängig prüfen. & Q12 \\
    \hline
    \textbf{P078} & N-Kanal-MOSFET & BSS89 & Typ. Selbstsperrender Kleinsignal-MOSFET; Pinbelegung herstellerabhängig prüfen. & Q12 \\
    \hline
    \textbf{P079} & Signalanschluss & Digitalausgang 0 & Aufdruck. Signalpegel und Belastbarkeit noch bestimmen. & -- \\
    \hline
    \textbf{P080} & Signalanschluss & Digitalausgang 1 & Aufdruck. Signalpegel und Belastbarkeit noch bestimmen. & -- \\
    \hline
    \textbf{P081} & Signalanschluss & Digitalausgang 2 & Aufdruck. Signalpegel und Belastbarkeit noch bestimmen. & -- \\
    \hline
    \textbf{P082} & Signalanschluss & Digitalausgang 3 & Aufdruck. Signalpegel und Belastbarkeit noch bestimmen. & -- \\
    \hline
    \textbf{P083} & Draht / Drahtbrücke & rot & Aufdruck. Leiterquerschnitt und Funktion nicht dokumentiert. & -- \\
    \hline
    \textbf{P084} & Widerstand & 120 k\ensuremath{\Omega}; \ensuremath{\pm}1 \% & Farbcode.  Ringe: Braun-Rot-Schwarz-Orange-Braun. & Q1 \\
    \hline
    \textbf{P085} & Leistungs-MOSFET, vermutlich & BUZ10 & vorläufig. „BUZ1U“ bzw. „BUZ 108048“ vermutlich BUZ10 mit Zusatzcode 8048. Vollständigen Aufdruck prüfen. & Q13 \\
    \hline
    \textbf{P086} & Widerstand & 10 k\ensuremath{\Omega}; \ensuremath{\pm}1 \% & Farbcode.  Ringe: Braun-Schwarz-Schwarz-Rot-Braun. & Q1 \\
    \hline
    \textbf{P087} & Widerstand & 150 \ensuremath{\Omega}; \ensuremath{\pm}5 \% & Farbcode.  Ringe: Braun-Grün-Braun-Gold. & Q1 \\
    \hline
    \textbf{P088} & Widerstand & 10 k\ensuremath{\Omega}; \ensuremath{\pm}5 \% & Farbcode.  Ringe: Braun-Schwarz-Orange-Gold. & Q1 \\
    \hline
    \textbf{P089} & Kondensator, vermutlich Keramik & vermutlich 22 nF & vorläufig. A5C223M: 223 = 22 000 pF = 22 nF; M üblicherweise \ensuremath{\pm}20 \%. Serie/Spannungsfestigkeit ungeklärt. & Q15 \\
    \hline
    \textbf{P090} & Widerstand & 10 k\ensuremath{\Omega}; \ensuremath{\pm}5 \% & Farbcode.  Ringe: Braun-Schwarz-Orange-Gold. & Q1 \\
    \hline
    \textbf{P091} & Kondensator, vermutlich Keramik & vermutlich 22 nF & vorläufig. A5C223M: 223 = 22 000 pF = 22 nF; M üblicherweise \ensuremath{\pm}20 \%. Serie/Spannungsfestigkeit ungeklärt. & Q15 \\
    \hline
    \textbf{P092} & Widerstand & 10 k\ensuremath{\Omega}; \ensuremath{\pm}1 \% & Farbcode.  Ringe: Braun-Schwarz-Schwarz-Rot-Braun. & Q1 \\
    \hline
    \textbf{P093} & Widerstand & 4,7 k\ensuremath{\Omega}; \ensuremath{\pm}1 \% & Farbcode.  Ringe: Gelb-Lila-Schwarz-Braun-Braun. & Q1 \\
    \hline
    \textbf{P094} & Widerstand & 390 k\ensuremath{\Omega}; \ensuremath{\pm}1 \% & Farbcode. 6. Ring rot: Zusatzkennzeichnung; Bedeutung/Temperaturkoeffizient prüfen. Ringe: Orange-Weiß-Schwraz-Orange-Braun-Rot. & Q1 \\
    \hline
    \textbf{P095} & Widerstand & 190 k\ensuremath{\Omega}; \ensuremath{\pm}1 \% & Farbcode.  Ringe: Braun-Weiß-Schwarz-Orange-Braun. & Q1 \\
    \hline
    \textbf{P096} & Trimmpotentiometer & Widerstand unbekannt & offen. 90P R5K 8235S nicht eindeutig entschlüsselt. 5 k\ensuremath{\Omega} ist nur eine Möglichkeit; Endanschlüsse messen. & -- \\
    \hline
    \textbf{P097} & Widerstand & 10 k\ensuremath{\Omega}; \ensuremath{\pm}5 \% & Farbcode.  Ringe: Braun-Schwarz-Orange-Gold. & Q1 \\
    \hline
    \textbf{P098} & Kondensator, vermutlich Keramik & vermutlich 22 nF & vorläufig. A5C223M: 223 = 22 000 pF = 22 nF; M üblicherweise \ensuremath{\pm}20 \%. Serie/Spannungsfestigkeit ungeklärt. & Q15 \\
    \hline
    \textbf{P099} & Widerstand & 10 k\ensuremath{\Omega}; \ensuremath{\pm}1 \% & Farbcode.  Ringe: Braun-Schwarz-Schwarz-Rot-Braun. & Q1 \\
    \hline
    \textbf{P100} & Widerstand & 10 k\ensuremath{\Omega}; \ensuremath{\pm}5 \% & Farbcode.  Ringe: Braun-Schwarz-Orange-Gold. & Q1 \\
    \hline
    \textbf{P101} & Kondensator, vermutlich Keramik & vermutlich 22 nF & vorläufig. A5C223M: 223 = 22 000 pF = 22 nF; M üblicherweise \ensuremath{\pm}20 \%. Serie/Spannungsfestigkeit ungeklärt. & Q15 \\
    \hline
    \textbf{P102} & Widerstand & 10 k\ensuremath{\Omega}; \ensuremath{\pm}1 \% & Farbcode.  Ringe: Braun-Schwarz-Schwarz-Rot-Braun. & Q1 \\
    \hline
    \textbf{P103} & Widerstand & 10 k\ensuremath{\Omega}; \ensuremath{\pm}1 \% & Farbcode.  Ringe: Braun-Schwrz-Schwarz-Rot-Braun. & Q1 \\
    \hline
    \textbf{P104} & Widerstand & 68 k\ensuremath{\Omega}; \ensuremath{\pm}1 \% & Farbcode.  Ringe: Blau-Grau-Schwarz-Rot-Braun. & Q1 \\
    \hline
    \textbf{P105} & Widerstand & 33 k\ensuremath{\Omega}; \ensuremath{\pm}1 \% & Farbcode.  Ringe: Orange-Orange-Schwarz-Rot-Braun. & Q1 \\
    \hline
    \textbf{P106} & Widerstand & 33,2 k\ensuremath{\Omega}; \ensuremath{\pm}1 \% & Farbcode.  Ringe: Orange-Orange-Rot-Rot-Braun. & Q1 \\
    \hline
    \textbf{P107} & PNP-Transistor & BC307B & Typ. Kleinsignaltransistor zum Schalten und Verstärken; konkrete Funktion aus Beschaltung ableiten. & Q11 \\
    \hline
    \textbf{P108} & Widerstand & 47,5 k\ensuremath{\Omega}; \ensuremath{\pm}1 \% & Farbcode.  Ringe: Gelb-Lila-Grün-Rot-Braun. & Q1 \\
    \hline
    \textbf{P109} & Widerstand & 33,2 k\ensuremath{\Omega}; \ensuremath{\pm}1 \% & Farbcode.  Ringe: Orange-Orange-Rot-Rot-Braun. & Q1 \\
    \hline
    \textbf{P110} & Diode, Typ ungeklärt & möglicherweise 1N4148 & vorläufig. „PN4841“ und ähnliche Lesarten könnten zum Aufdruck 4148PH gehören. Nicht als Z-Diode bestätigt. Aufdruck/Foto prüfen. & Q10 \\
    \hline
    \textbf{P111} & Widerstand & 2,74 k\ensuremath{\Omega}; \ensuremath{\pm}1 \% & Farbcode.  Ringe: Rot-Lila-Gelb-Braun-Braun. & Q1 \\
    \hline
    \textbf{P112} & Widerstand & 100 k\ensuremath{\Omega}; \ensuremath{\pm}1 \% & Farbcode.  Ringe: Braun-Schwarz-Schwarz-Orange-Braun. & Q1 \\
    \hline
    \textbf{P113} & Trimmpotentiometer & Widerstand unbekannt & offen. 7633D kann ein Fertigungscode sein; daraus keinen Widerstandswert ableiten. & -- \\
    \hline
    \textbf{P114} & Widerstand & 1 M\ensuremath{\Omega}; \ensuremath{\pm}1 \% & Farbcode.  Ringe: Braun-Schwarz-Schwarz-Gelb-Braun. & Q1 \\
    \hline
    \textbf{P115} & Widerstand & 10 k\ensuremath{\Omega}; \ensuremath{\pm}1 \% & Farbcode.  Ringe: Braun-Schwarz-Schwarz-Rot-Braun. & Q1 \\
    \hline
    \textbf{P116} & Widerstand & 100 k\ensuremath{\Omega}; \ensuremath{\pm}1 \% & Farbcode.  Ringe: Braun-Schwarz-Schwarz-Orange-Braun. & Q1 \\
    \hline
    \textbf{P117} & Widerstand & 10 k\ensuremath{\Omega}; \ensuremath{\pm}1 \% & Farbcode.  Ringe: Braun-Schwarz-Schwarz-Rot-braun. & Q1 \\
    \hline
    \textbf{P118} & Folienkondensator & 1 \ensuremath{\mu}F / vermutlich 63 V & Aufdruck. WIMA, Aufdruck 1uF / 63-N / V3. Serie und Spannungsart prüfen. & -- \\
    \hline
    \textbf{P119} & Diode, Typ ungeklärt & möglicherweise 1N4148 & vorläufig. „PN4841“ und ähnliche Lesarten könnten zum Aufdruck 4148PH gehören. Nicht als Z-Diode bestätigt. Aufdruck/Foto prüfen. & Q10 \\
    \hline
    \textbf{P120} & Widerstand & 10 M\ensuremath{\Omega} & Aufdruck. Toleranz und Belastbarkeit nicht dokumentiert. & -- \\
    \hline
    \textbf{P121} & Widerstand & 1 M\ensuremath{\Omega}; \ensuremath{\pm}1 \% & Farbcode.  Ringe: Braun-Schwarz-Schwarz-Gelb-Braun. & Q1 \\
    \hline
    \textbf{P122} & Potentiometer, vermutlich & vermutlich 20 k\ensuremath{\Omega} & vorläufig. Aufdruck M20K L8; 20K als Widerstand interpretiert. Kennlinie und Gesamtwiderstand prüfen. & -- \\
    \hline
    \textbf{P123} & Bauteil, nicht identifiziert & 565-2 / 0,4 A & offen. 0,4 A ist notiert, die Bauteilart bleibt ungeklärt. Foto und vollständige Kennzeichnung erforderlich. & -- \\
    \hline
    \textbf{P124} & Widerstand & 47 k\ensuremath{\Omega}; \ensuremath{\pm}1 \% & Farbcode.  Ringe: Gelb-Lila-Schwarz-Rot-Braun. & Q1 \\
    \hline
    \textbf{P125} & Diode, Typ ungeklärt & möglicherweise 1N4148 & vorläufig. „PN4841“ und ähnliche Lesarten könnten zum Aufdruck 4148PH gehören. Nicht als Z-Diode bestätigt. Aufdruck/Foto prüfen. & Q10 \\
    \hline
    \textbf{P126} & Kondensator, vermutlich & möglicherweise 100 nF & vorläufig. Aufdruck 1042 / Mexiko: nur falls 104 der Kapazitätscode ist, 100 nF. Vierstelliger Gesamtcode nicht eindeutig. & Q15 \\
    \hline
    \textbf{P127} & Widerstand & 680 k\ensuremath{\Omega}; \ensuremath{\pm}1 \% & Farbcode. 6. Ring rot: Zusatzkennzeichnung; Bedeutung/Temperaturkoeffizient prüfen. Ringe: Blau-Grau-Schwarz-Orange-Braun-Rot. & Q1 \\
    \hline
    \textbf{P128} & Potentiometer & 10 k\ensuremath{\Omega} & Aufdruck. Aufdruck 10kO 842; Kennlinie und Bauform ungeklärt. & -- \\
    \hline
    \textbf{P129} & Widerstand & 390 k\ensuremath{\Omega}; \ensuremath{\pm}1 \% & Farbcode. 6. Ring rot: Zusatzkennzeichnung; Bedeutung/Temperaturkoeffizient prüfen. Ringe: Orange-Weiß-Schwarz-Orange-Braun-Rot. & Q1 \\
    \hline
    \textbf{P130} & Widerstand & 9,09 k\ensuremath{\Omega}; \ensuremath{\pm}1 \% & Farbcode.  Ringe: Weiß-Schwarz-Weiß-Braun-Braun. & Q1 \\
    \hline
    \textbf{P131} & Widerstand & 9,09 k\ensuremath{\Omega}; \ensuremath{\pm}1 \% & Farbcode.  Ringe: Weiß-Schwarz-Weiß-Braun-Braun. & Q1 \\
    \hline
    \textbf{P132} & Trimmpotentiometer & 5 k\ensuremath{\Omega} & Aufdruck. Aufdruck 5kO 7526D; in der Aufnahme als möglicherweise defekt vermerkt. Funktionsprüfung erforderlich. & -- \\
    \hline
    \textbf{P133} & Widerstand & 47 k\ensuremath{\Omega}; \ensuremath{\pm}1 \% & Farbcode.  Ringe: Gelb-Lila-Schwarz-Rot-Braun. & Q1 \\
    \hline
    \textbf{P134} & Kondensator, vermutlich & möglicherweise 100 nF & vorläufig. Aufdruck 1042 / Mexiko: nur falls 104 der Kapazitätscode ist, 100 nF. Vierstelliger Gesamtcode nicht eindeutig. & Q15 \\
    \hline
    \textbf{P135} & Widerstand & 1 M\ensuremath{\Omega}; \ensuremath{\pm}1 \% & Farbcode.  Ringe: Braun-Schwarz-Schwarz-Gelb-Braun. & Q1 \\
    \hline
    \textbf{P136} & Widerstand & 100 k\ensuremath{\Omega}; \ensuremath{\pm}1 \% & Farbcode.  Ringe: Braun-Schwarz-Schwarz-Orange-Braun. & Q1 \\
    \hline
    \textbf{P137} & Widerstand & 100 k\ensuremath{\Omega}; \ensuremath{\pm}1 \% & Farbcode.  Ringe: Braun-Schwarz-Schwarz-Orange-Braun. & Q1 \\
    \hline
    \textbf{P138} & Diode, Typ ungeklärt & möglicherweise 1N4148 & vorläufig. „PN4841“ und ähnliche Lesarten könnten zum Aufdruck 4148PH gehören. Nicht als Z-Diode bestätigt. Aufdruck/Foto prüfen. & Q10 \\
    \hline
    \textbf{P139} & Widerstand & 1 M\ensuremath{\Omega}; \ensuremath{\pm}1 \% & Farbcode.  Ringe: Braun-Schwarz-Schwarz-Gelb-Braun. & Q1 \\
    \hline
    \textbf{P140} & Widerstand & 1 M\ensuremath{\Omega}; \ensuremath{\pm}1 \% & Farbcode.  Ringe: Braun-Schwarz-Schwarz-Gelb-Braun. & Q1 \\
    \hline
    \textbf{P141} & Widerstand & 1 M\ensuremath{\Omega}; \ensuremath{\pm}1 \% & Farbcode.  Ringe: Braun-Schwarz-Schwarz-Gelb-Braun. & Q1 \\
    \hline
    \textbf{P142} & Widerstand & 220 k\ensuremath{\Omega}; \ensuremath{\pm}1 \% & Farbcode. 6. Ring rot: Zusatzkennzeichnung; Bedeutung/Temperaturkoeffizient prüfen. Ringe: Rot-Rot-Schwarz-Orange-Braun-Rot. & Q1 \\
    \hline
    \textbf{P143} & Widerstand & 220 k\ensuremath{\Omega}; \ensuremath{\pm}1 \% & Farbcode. 6. Ring rot: Zusatzkennzeichnung; Bedeutung/Temperaturkoeffizient prüfen. Ringe: Rot-Rot-Schwarz-Orange-Braun-Rot. & Q1 \\
    \hline
    \textbf{P144} & Widerstand & 220 k\ensuremath{\Omega}; \ensuremath{\pm}1 \% & Farbcode. 6. Ring rot: Zusatzkennzeichnung; Bedeutung/Temperaturkoeffizient prüfen. Ringe: Rot-Rot-Schwarz-Orange-Braun-Rot. & Q1 \\
    \hline
    \textbf{P145} & Widerstand & 220 k\ensuremath{\Omega}; \ensuremath{\pm}1 \% & Farbcode. 6. Ring rot: Zusatzkennzeichnung; Bedeutung/Temperaturkoeffizient prüfen. Ringe: Rot-Rot-Schwarz-Orange-Braun-Rot. & Q1 \\
    \hline
    \textbf{P146} & Widerstand & 100 k\ensuremath{\Omega}; \ensuremath{\pm}1 \% & Farbcode.  Ringe: Braun-Schwarz-Schwarz-Orage-Braun. & Q1 \\
    \hline
    \textbf{P147} & Diode, Typ ungeklärt & unbekannt & offen. Als Z-Diode erfasst, jedoch ohne Aufdruck. Bauteilart und gegebenenfalls Zenerspannung bestimmen. & -- \\
    \hline
    \textbf{P148} & Widerstand & 100 k\ensuremath{\Omega}; \ensuremath{\pm}1 \% & Farbcode.  Ringe: Braun-Schwarz-Schwarz-Orage-Braun. & Q1 \\
    \hline
    \textbf{P149} & Diode, Typ ungeklärt & möglicherweise 1N4148 & vorläufig. „PN4841“ und ähnliche Lesarten könnten zum Aufdruck 4148PH gehören. Nicht als Z-Diode bestätigt. Aufdruck/Foto prüfen. & Q10 \\
    \hline
    \textbf{P150} & Widerstand & 100 k\ensuremath{\Omega}; \ensuremath{\pm}1 \% & Farbcode.  Ringe: Braun-Schwarz-Schwarz-Orage-Braun. & Q1 \\
    \hline
    \textbf{P151} & Diode, Typ ungeklärt & möglicherweise 1N4148 & vorläufig. „PN4841“ und ähnliche Lesarten könnten zum Aufdruck 4148PH gehören. Nicht als Z-Diode bestätigt. Aufdruck/Foto prüfen. & Q10 \\
    \hline
    \textbf{P152} & Kondensator, vermutlich & möglicherweise 100 nF & vorläufig. Aufdruck 1042 / Mexiko: nur falls 104 der Kapazitätscode ist, 100 nF. Vierstelliger Gesamtcode nicht eindeutig. & Q15 \\
    \hline
    \textbf{P153} & Kondensator, vermutlich & möglicherweise 100 nF & vorläufig. Aufdruck 1042 / Mexiko: nur falls 104 der Kapazitätscode ist, 100 nF. Vierstelliger Gesamtcode nicht eindeutig. & Q15 \\
    \hline
    \textbf{P154} & Kondensator, vermutlich & möglicherweise 100 nF & vorläufig. Aufdruck 1042 / Mexiko: nur falls 104 der Kapazitätscode ist, 100 nF. Vierstelliger Gesamtcode nicht eindeutig. & Q15 \\
    \hline
    \textbf{P155} & Kondensator, vermutlich & möglicherweise 100 nF & vorläufig. Aufdruck 1042 / Mexiko: nur falls 104 der Kapazitätscode ist, 100 nF. Vierstelliger Gesamtcode nicht eindeutig. & Q15 \\
    \hline
    \textbf{P156} & Kondensator, vermutlich & möglicherweise 100 nF & vorläufig. Aufdruck 1042 / Mexiko: nur falls 104 der Kapazitätscode ist, 100 nF. Vierstelliger Gesamtcode nicht eindeutig. & Q15 \\
    \hline
    \textbf{P157} & Kondensator, vermutlich & möglicherweise 100 nF & vorläufig. Aufdruck 1042 / Mexiko: nur falls 104 der Kapazitätscode ist, 100 nF. Vierstelliger Gesamtcode nicht eindeutig. & Q15 \\
    \hline
    \textbf{P158} & Kondensator, vermutlich & möglicherweise 100 nF & vorläufig. Aufdruck 1042 / Mexiko: nur falls 104 der Kapazitätscode ist, 100 nF. Vierstelliger Gesamtcode nicht eindeutig. & Q15 \\
    \hline
    \textbf{P159} & Widerstand & 100 k\ensuremath{\Omega}; \ensuremath{\pm}1 \% & Farbcode.  Ringe: Braun-Schwarz-Schwarz-Orange-Braun. & Q1 \\
    \hline
    \textbf{P160} & Widerstand & 100 k\ensuremath{\Omega}; \ensuremath{\pm}1 \% & Farbcode.  Ringe: Braun-Schwarz-Schwarz-Orange-Braun. & Q1 \\
    \hline
    \textbf{P161} & Widerstand & 100 k\ensuremath{\Omega}; \ensuremath{\pm}1 \% & Farbcode.  Ringe: Braun-Schwarz-Schwarz-Orange-Braun. & Q1 \\
    \hline
    \textbf{P162} & Trimmpotentiometer & Widerstand unbekannt & offen. Aufdruck 100 821 nicht eindeutig: 100 kann je nach Kennzeichnung 100 \ensuremath{\Omega} oder ein Code sein. Messen. & -- \\
    \hline
    \textbf{P163} & Widerstand & 82 k\ensuremath{\Omega}; \ensuremath{\pm}1 \% & Farbcode.  Ringe: Grau-Rot-Schwarz-Rot-Braun. & Q1 \\
    \hline
    \textbf{P164} & Widerstand & 4,75 k\ensuremath{\Omega}; \ensuremath{\pm}1 \% & Farbcode.  Ringe: Gelb-Lila-Grün-Braun-Braun. & Q1 \\
    \hline
    \textbf{P165} & PNP-Transistor & BC307B & Typ. Kleinsignaltransistor zum Schalten und Verstärken; konkrete Funktion aus Beschaltung ableiten. & Q11 \\
    \hline
    \textbf{P166} & Widerstand & 470 k\ensuremath{\Omega}; \ensuremath{\pm}1 \% & Farbcode.  Ringe: Gelb-Lila-Schwarz-Orange-Braun. & Q1 \\
    \hline
    \textbf{P167} & Spannungsreferenz & LM329DZ / ca. 6,9 V & Typ. Kein diskreter Transistor. „524“ ist Zusatzaufdruck, kein Referenzspannungswert. & Q3 \\
    \hline
    \textbf{P168} & Widerstand & 2 k\ensuremath{\Omega}; \ensuremath{\pm}5 \% & Farbcode.  Ringe: Rot-Schwarz-Rot-Gold. & Q1 \\
    \hline
\end{GPMTable}

\subsubsection*{Quellen zur Bauteilzuordnung}
\begin{description}[style=nextline,leftmargin=1.2cm]
\item[Q1] Vishay: Widerstandsfarbcode. \url{https://www.vishay.com/docs/49411/resistor_color_code_calculator.pdf} (Abruf: 27.10.2026).
\item[Q2] Texas Instruments: TL497A. \url{https://www.ti.com/product/TL497A} (Abruf: 27.10.2026).
\item[Q3] Analog Devices: LM329, 6,9-V-Spannungsreferenz. \url{https://www.analog.com/en/products/lm329.html} (Abruf: 27.10.2026).
\item[Q4] Texas Instruments: CD4001B, NOR-Gatter. \url{https://www.ti.com/product/CD4001B/part-details/CD4001BE} (Abruf: 27.10.2026).
\item[Q5] Renesas: ICL7621, zweifacher CMOS-Operationsverstärker. \url{https://www.renesas.com/en/products/icl7621} (Abruf: 27.10.2026).
\item[Q6] Nexperia: HEF4528B, zweifaches Monoflop. \url{https://www.nexperia.com/products/analog-logic-ics/logic/specialty-logic/multivibrators/serie/hef4528b} (Abruf: 27.10.2026).
\item[Q7] STMicroelectronics: HCF4051, Analogmultiplexer. \url{https://www.st.com/resource/en/datasheet/hcf4051.pdf} (Abruf: 27.10.2026).
\item[Q8] Texas Instruments: TL061. \url{https://www.ti.com/product/TL061} (Abruf: 27.10.2026).
\item[Q9] Vishay: BAT41, Schottkydiode. \url{https://www.vishay.com/docs/85659/bat41.pdf} (Abruf: 27.10.2026).
\item[Q10] NXP/Philips: 1N4148, Originaldatenblatt mit Aufdruck 4148PH. \url{https://assets.nexperia.com/documents/data-sheet/1N4148_1N4448.pdf} (Abruf: 27.10.2026).
\item[Q11] Philips: BC307/BC307B, Originaldatenblatt im Archiv. \url{https://www.z80ne.com/datasheets/bc307b.pdf} (Abruf: 27.10.2026).
\item[Q12] Siemens: SIPMOS-Datenbuch 1985, Druckseiten 7, 55 und 60 (BSS89/BSS91). \url{https://bitsavers.trailing-edge.com/components/siemens/1985_Siemens_SIPMOS_Small_Signal_Transistors__Power_Transistors.pdf} (Abruf: 27.10.2026).
\item[Q13] Siemens: SIPMOS-Datenbuch 1987/88, Druckseite 160 (BUZ10). \url{https://bitsavers.trailing-edge.com/components/siemens/1987_Siemens_SIPMOS_Components.pdf} (Abruf: 27.10.2026).
\item[Q14] Bourns: 4600X, dreistelliger Widerstandscode; kein Nachweis der konkreten Netzwerkserie. \url{https://www.bourns.com/docs/product-datasheets/4600x.pdf} (Abruf: 27.10.2026).
\item[Q15] KYOCERA AVX: Kapazitätscode und Toleranz M; kein Nachweis der konkreten Kondensatorserie. \url{https://catalogs.kyocera-avx.com/SurfaceMount.pdf} (Abruf: 27.10.2026).
\item[Q16] STMicroelectronics: L78L, positive 100-mA-Spannungsregler. \url{https://www.st.com/en/power-management/l78l.html} (Abruf: 27.10.2026).
\item[Q17] Siemens: TAE1453, Datenblattbeschreibung im Archiv. \url{https://www.alldatasheet.com/datasheet-pdf/pdf/45790/SIEMENS/TAE1453.html} (Abruf: 27.10.2026).
\end{description}
